% GENERATED FILE. Edit canonical lessons, then run npm run generate:books.
% canonical-source-hash: fnv1a64:c29585c938e0bb7e
% canonical-lessons: SA-C289-abhijanati, SA-C289-varnayati, SA-C289-niriksate, SA-C289-cihnam, SA-C289-pratyaksam

\chapter{What the Witness Saw}
\label{ch:sa-what-the-witness-saw}
\hlchaptermodality{289}

\begin{chapteropening}
\textbf{By the end of this chapter:} \emph{I can say what I saw and what would have been different.}

Complete the last lesson of chapter 289: I can say what I saw and what would have been different.
\end{chapteropening}

\section[\texorpdfstring{\sk{अभिजानाति} (abhijānāti)}{अभिजानाति (abhijānāti)}]{\sk{अभिजानाति} (abhijānāti) — recognizes}

\label{lesson:SA-C289-abhijanati}



\begin{quote}
Before the new one: say the Sanskrit for threatens, then the Sanskrit for a weapon.
\end{quote}

\subsection*{You'll want to know: \sk{अभिजानाति}}
\textbf{\sk{अभिजानाति}} — \emph{abhijānāti} — \textquotedblleft{}recognizes\textquotedblright{}.

\begin{quote}
\sk{यदि} \sk{दीपः} \sk{अभविष्यत्}, \sk{साक्षी} \sk{चौरम्} \sk{अभ्यज्ञास्यत्।}
\end{quote}

\begin{quote}
\emph{yadi dīpaḥ abhaviṣyat, sākṣī cauram abhyajñāsyat.}
\end{quote}

\begin{quote}
If there had been a lamp, the witness would have recognized the thief.
\end{quote}

Past (he …): \sk{सः} \sk{चौरम्} \sk{अभिज्ञातवान्}, \emph{saḥ cauram abhijñātavān}.

Future: \sk{अभिज्ञास्यति}, \emph{abhijñāsyati}.

\begin{cousinweb}
\textbf{\sk{अभिजानाति}} (\emph{abhijānāti}) is built from \emph{abhi}, \textquotedblleft{}towards\textquotedblright{}, and the root √jñā, \textquotedblleft{}to know\textquotedblright{}; MW gives \emph{abhi-jñā} as to recognize or become aware of (MW).
\end{cousinweb}

\begin{grammarlens}[title={What you've built}]
One more word for telling what happened.
\end{grammarlens}

\subsection*{Your turn}
\begin{itemize}
\raggedright
  \item \emph{Say it:} \emph{abhijānāti}
  \item \emph{Say it:} \emph{abhijānāti}, once more
  \item \emph{Say it:} \emph{yadi dīpaḥ abhaviṣyat, sākṣī cauram abhyajñāsyat.} — If there had been a lamp, the witness would have recognized the thief.
  \item \emph{Recall:} say the Sanskrit for threatens, then the Sanskrit for a weapon, then say \emph{abhijānāti} again
\end{itemize}

\begin{tcolorbox}[breakable,colback=teal!4,colframe=teal!35!black,title={Before you move on}]
What does \sk{अभिजानाति} mean? (\textquotedblleft{}Recognizes\textquotedblright{}.) Say it once more.
\end{tcolorbox}
\section[\texorpdfstring{\sk{वर्णयति} (varṇayati)}{वर्णयति (varṇayati)}]{\sk{वर्णयति} (varṇayati) — describes}

\label{lesson:SA-C289-varnayati}



\begin{quote}
Before the new one: say the Sanskrit for a weapon, then the Sanskrit for recognizes.
\end{quote}

\subsection*{You'll want to know: \sk{वर्णयति}}
\textbf{\sk{वर्णयति}} — \emph{varṇayati} — \textquotedblleft{}describes\textquotedblright{}.

\begin{quote}
\sk{साक्षी} \sk{चौरस्य} \sk{मुखं} \sk{वर्णयति।}
\end{quote}

\begin{quote}
\emph{sākṣī caurasya mukhaṃ varṇayati.}
\end{quote}

\begin{quote}
The witness describes the thief's face.
\end{quote}

Past (he …): \sk{सः} \sk{मुखं} \sk{वर्णितवान्}, \emph{saḥ mukhaṃ varṇitavān}.

\begin{cousinweb}
\textbf{\sk{वर्णयति}} (\emph{varṇayati}) is built on \emph{varṇa}, \textquotedblleft{}colour, outward appearance\textquotedblright{}; MW treats \emph{varṇayati} as a verb made from that noun, to paint and to describe (MW).
\end{cousinweb}

\begin{grammarlens}[title={What you've built}]
One more word for telling what happened.
\end{grammarlens}

\subsection*{Your turn}
\begin{itemize}
\raggedright
  \item \emph{Say it:} \emph{varṇayati}
  \item \emph{Say it:} \emph{varṇayati}, once more
  \item \emph{Say it:} \emph{sākṣī caurasya mukhaṃ varṇayati.} — The witness describes the thief's face.
  \item \emph{Recall:} say the Sanskrit for a weapon, then the Sanskrit for recognizes, then say \emph{varṇayati} again
\end{itemize}

\begin{tcolorbox}[breakable,colback=teal!4,colframe=teal!35!black,title={Before you move on}]
What does \sk{वर्णयति} mean? (\textquotedblleft{}Describes\textquotedblright{}.) Say it once more.
\end{tcolorbox}
\section[\texorpdfstring{\sk{निरीक्षते} (nirīkṣate)}{निरीक्षते (nirīkṣate)}]{\sk{निरीक्षते} (nirīkṣate) — observes closely}

\label{lesson:SA-C289-niriksate}



\begin{quote}
Before the new one: say the Sanskrit for recognizes, then the Sanskrit for describes.
\end{quote}

\subsection*{You'll want to know: \sk{निरीक्षते}}
\textbf{\sk{निरीक्षते}} — \emph{nirīkṣate} — \textquotedblleft{}observes closely\textquotedblright{}.

\begin{quote}
\sk{साक्षी} \sk{सर्वं} \sk{निरीक्षते।}
\end{quote}

\begin{quote}
\emph{sākṣī sarvaṃ nirīkṣate.}
\end{quote}

\begin{quote}
The witness observes everything closely.
\end{quote}

Past (he …): \sk{सः} \sk{सर्वं} \sk{निरीक्षितवान्}, \emph{saḥ sarvaṃ nirīkṣitavān}.

\begin{cousinweb}
\textbf{\sk{निरीक्षते}} (\emph{nirīkṣate}) is built from \emph{nir} and the root √īkṣ, \textquotedblleft{}to see\textquotedblright{}; MW gives \emph{nir-īkṣ} as to look at, regard and observe (MW).
\end{cousinweb}

\begin{grammarlens}[title={What you've built}]
One more word for telling what happened.
\end{grammarlens}

\subsection*{Your turn}
\begin{itemize}
\raggedright
  \item \emph{Say it:} \emph{nirīkṣate}
  \item \emph{Say it:} \emph{nirīkṣate}, once more
  \item \emph{Say it:} \emph{sākṣī sarvaṃ nirīkṣate.} — The witness observes everything closely.
  \item \emph{Recall:} say the Sanskrit for recognizes, then the Sanskrit for describes, then say \emph{nirīkṣate} again
\end{itemize}

\begin{tcolorbox}[breakable,colback=teal!4,colframe=teal!35!black,title={Before you move on}]
What does \sk{निरीक्षते} mean? (\textquotedblleft{}Observes closely\textquotedblright{}.) Say it once more.
\end{tcolorbox}
\section[\texorpdfstring{\sk{चिह्नम्} (cihnam)}{चिह्नम् (cihnam)}]{\sk{चिह्नम्} (cihnam) — a trace; a sign}

\label{lesson:SA-C289-cihnam}



\begin{quote}
Before the new one: say the Sanskrit for describes, then the Sanskrit for observes closely.
\end{quote}

\subsection*{You'll want to know: \sk{चिह्नम्}}
\textbf{\sk{चिह्नम्}} — \emph{cihnam} — \textquotedblleft{}a trace; a sign\textquotedblright{}.

\begin{quote}
\sk{भित्तौ} \sk{चौरस्य} \sk{चिह्नम्} \sk{अस्ति।}
\end{quote}

\begin{quote}
\emph{bhittau caurasya cihnam asti.}
\end{quote}

\begin{quote}
There is a mark of the thief on the wall.
\end{quote}

A neuter noun.

\begin{cousinweb}
\textbf{\sk{चिह्नम्}} (\emph{cihnam}) MW gives \emph{cihna} as a mark, a sign or a characteristic, without a derivation, so its origin is uncertain (MW).
\end{cousinweb}

\begin{grammarlens}[title={What you've built}]
One more word for telling what happened.
\end{grammarlens}

\subsection*{Your turn}
\begin{itemize}
\raggedright
  \item \emph{Say it:} \emph{cihnam}
  \item \emph{Say it:} \emph{cihnam}, once more
  \item \emph{Say it:} \emph{bhittau caurasya cihnam asti.} — There is a mark of the thief on the wall.
  \item \emph{Recall:} say the Sanskrit for describes, then the Sanskrit for observes closely, then say \emph{cihnam} again
\end{itemize}

\begin{tcolorbox}[breakable,colback=teal!4,colframe=teal!35!black,title={Before you move on}]
What does \sk{चिह्नम्} mean? (\textquotedblleft{}A trace; a sign\textquotedblright{}.) Say it once more.
\end{tcolorbox}
\section[\texorpdfstring{\sk{प्रत्यक्षम्} (pratyakṣam)}{प्रत्यक्षम् (pratyakṣam)}]{\sk{प्रत्यक्षम्} (pratyakṣam) — before one's eyes; in person}

\label{lesson:SA-C289-pratyaksam}



\begin{quote}
Before the new one: say the Sanskrit for observes closely, then the Sanskrit for a trace.
\end{quote}

\subsection*{You'll want to know: \sk{प्रत्यक्षम्}}
\textbf{\sk{प्रत्यक्षम्}} — \emph{pratyakṣam} — \textquotedblleft{}before one's eyes; in person\textquotedblright{}.

\begin{quote}
\sk{अहं} \sk{प्रत्यक्षं} \sk{दृष्टवान्।}
\end{quote}

\begin{quote}
\emph{ahaṃ pratyakṣaṃ dṛṣṭavān.}
\end{quote}

\begin{quote}
I saw it with my own eyes.
\end{quote}

\begin{cousinweb}
\textbf{\sk{प्रत्यक्षम्}} (\emph{pratyakṣam}) is built from \emph{prati}, \textquotedblleft{}before, facing\textquotedblright{}, and \emph{akṣa}, \textquotedblleft{}the eye\textquotedblright{}; MW gives \emph{pratyakṣa} as present before the eyes, visible (MW).
\end{cousinweb}

\begin{grammarlens}[title={What you've built}]
One more word for telling what happened.
\end{grammarlens}

\subsection*{Your turn}
\begin{itemize}
\raggedright
  \item \emph{Say it:} \emph{pratyakṣam}
  \item \emph{Say it:} \emph{pratyakṣam}, once more
  \item \emph{Say it:} \emph{ahaṃ pratyakṣaṃ dṛṣṭavān.} — I saw it with my own eyes.
  \item \emph{Recall:} say the Sanskrit for observes closely, then the Sanskrit for a trace, then say \emph{pratyakṣam} again
\end{itemize}

\begin{tcolorbox}[breakable,colback=teal!4,colframe=teal!35!black,title={Before you move on}]
What does \sk{प्रत्यक्षम्} mean? (\textquotedblleft{}Before one's eyes; in person\textquotedblright{}.) Say it once more.
\end{tcolorbox}
